% Appendix B
\chapter{Reproducibility and Sample Code}\label{app:code}

The implementation was written for this thesis and does not import code from
other saliency projects. This appendix lists its main parts, the points at
which the work departed from its own procedure, the commands that produce the
main results, and the forward pass of the decoder.

\section{Organisation of the Code}

\Cref{tab:code} lists the main files. The research log records every decision
rule with its date, written before the experiment it governs, and every
result with the folder it was read from.

\begin{table}[htbp]
  \centering
  \caption{Main files of the implementation}\label{tab:code}
  \footnotesize
  \begin{tabularx}{\textwidth}{@{}>{\ttfamily\raggedright}p{6.2cm}>{\raggedright\arraybackslash}X@{}}
    \toprule
    \textrm{\textbf{File}} & \textbf{Content} \\
    \midrule
    surge/backbones.py & frozen encoders (\acrshort{V-JEPA}~2 \acrshort{ViT}-L and \acrshort{ViT}-g, \acrshort{V-JEPA}~2.1, \acrshort{UMT}-L, DINOv2-L) and their taps; predictor error for gate G2 and the batched masked-tube readout of the predictability gate \\
    surge/model.py & feature pyramid, (2+1)D decoder, auxiliary heads, centre prior, surprise and audio modules, second-encoder fusion, stream buffer of the decoder state \\
    surge/data.py & dataset readers, clip sampling, pairs of adjacent windows, targets, inference windows \\
    surge/losses.py & training loss (\cref{eq:loss}) \\
    surge/metrics.py & \acrshort{CC}, \acrshort{NSS}, \acrshort{SIM}, \acrshort{KLD}, \acrshort{AUC-J} and \acrshort{sAUC} \\
    surge/infer.py & whole-video inference (\cref{alg:inference}) and scoring \\
    surge/lora.py & low-rank adaptation (\cref{eq:lora}) \\
    scripts/train.py & training with holdout selection and clean stop/resume \\
    scripts/evaluate.py & scoring of models, ensembles or saved maps \\
    scripts/tune\_sigma.py & selection of the output blur on the holdout \\
    scripts/compare.py & paired bootstrap over videos \\
    scripts/g2\_surprise.py & gate G2 \\
    scripts/prep\_hollywood\_gaze.py & Hollywood-2 (Vig gaze) maps \\
    scripts/prep\_stavis\_audio.py & audio features for the audio experiment \\
    scripts/thesis\_qualitative.py & qualitative examples and cost measurements \\
    scripts/g2\_intervals.py & bootstrap intervals of gate G2 \\
    scripts/ablation\_seeds.py & component ablations with three seeds \\
    scripts/review\_tables.py & numbers of the checks added after the first draft \\
    scripts/g2\_maps.py & prediction-error maps with both mask tokens, the masked-tube readout and the checks of the visible-token encoder (\cref{sec:surprise}) \\
    scripts/g2\_incremental.py & incremental information (\cref{eq:incremental}) and the re-scored gate G2 \\
    scripts/make\_lc\_splits.py, scripts/learning\_curve.py & training subsets and analysis of the learning curve (\cref{sec:learning-curve}) \\
    scripts/time\_inference.py & measured inference time of SalFoM and SURGE-D (\cref{sec:cost}) \\
    scripts/count\_umt\_flops.py & operation count of SalFoM's encoder (\cref{sec:cost}) \\
    scripts/queue\_review\_a.sh to \_e.sh & the runs of the checks added after the first draft, in order \\
    scripts/asd\_speaker.py & feasibility check of the speaker maps (\cref{sec:speaker-maps}) \\
    scripts/ssl\_maps.py & sound-source localisation maps and the choice of the audio sampling rate (\cref{sec:ssl-maps}) \\
    scripts/summarize\_ssl.py & outcome of the sound-source maps under their rule \\
    scripts/queue\_14a.sh, \_14a2.sh, \_14d.sh, \_14e.sh, \_14f.sh & the runs of the extensions (\cref{sec:extensions}), in order \\
    scripts/train\_multi.py & several decoders trained together on one stream of encoder features; pairs of adjacent windows for the decoder state (\cref{sec:state}) \\
    scripts/probe\_eval.py & inference with windows that do not overlap, with or without the buffer, and accuracy by position in the window (\cref{fig:state}) \\
    scripts/inv\_summary.py, scripts/p1\_first\_window.py & tests of the decoder state under its rule; the analysis without the first window of each video \\
    scripts/queue\_16p1.sh & the runs of the decoder state, in order \\
    configs/ & configurations and all split and holdout lists \\
    research/log.md & dated decision rules and results \\
    \bottomrule
  \end{tabularx}
\end{table}

\section{Corrections and Deviations}\label{sec:deviations}

The procedure of \cref{sec:design} requires every rule to be written before
its experiment and to be applied unchanged. This section collects the points
at which the work departed from that procedure, or at which an error in an
earlier step was corrected, so that they need not be searched for in the log.

\paragraph{The mask token of the predictor.} Gate G2 was first run with the
mask token that the model-loading library selects by default, a zero vector
that pre-training never updated (\cref{sec:surprise}). The error was found in
the second follow-up, while the masked-tube readout was being prepared. Before
any score of the corrected maps against gaze data existed, the log recorded
that the trained token is the primary readout and what the thesis would report
for either outcome of the re-scored gate. The re-scored gate failed as its
first run had (\cref{sec:g2}), so the decision not to train the surprise
module stands. The numbers of the first run are kept in
\cref{tab:g2,tab:app-g2} as its record. The diagnostic on eleven clips without gaze data (\cref{sec:g2}) used
the untrained token and was not repeated. SURGE-D uses only the encoder, so no
model, score or comparison of SURGE-D is affected.

\paragraph{A test of the analysis programs.} After those decisions had been
recorded, the programs that compute the prediction-error maps and their
incremental information were tested end to end on four of the 60 holdout
videos. The test printed scores of these four videos against gaze data,
including scores with the trained token, and they were seen before the analysis
of all 60 videos had finished. No rule, margin, primary comparison or line of
the analysis programs was changed afterwards, and the four videos are part of
the full analysis.

\paragraph{Rules and analyses added after results.} The tie margin of 0.01
for comparisons with published numbers was fixed after the first such
comparison had been seen (\cref{sec:design}). The bootstrap intervals of the
first run of gate G2 were specified and computed after its outcome was known.
The comparison of each ensemble with its seed-0 member (\cref{tab:app-ens}) was
added after the scores that it uses. None of the three decided a model or a
setting.

\paragraph{The operation count of SalFoM's encoder.} The first draft gave 1.90
\acrshort{TFLOPs} for SalFoM's encoder. That count omitted the products with the
attention matrices, which the counts of SURGE-D include; with them, it is 2.87
\acrshort{TFLOPs} (\cref{sec:cost}). The correction raises SalFoM's encoder
computation per output frame from about 17 to about 25 times that of SURGE-D
with \acrshort{ViT}-L and from about 6 to about 9 times that with
\acrshort{ViT}-g, and it makes a causal \acrshort{ViT}-L model cheaper than
SalFoM's encoder rather than comparable to it.

\paragraph{The measured times.} The protocol of the timing required that
nothing else be started on the machine while it ran. Short commands were
started during the first run, and two brief file searches during the second;
neither used the \acrshort{GPU}. The measurement was therefore made a third time
under the same rule, which it met, and the third run is reported
(\cref{tab:timing}). The criterion was applied to all three runs, and their
medians differ by at most 5.1\,\%.

\paragraph{The second look at the second encoder.} The rule of the second
encoder provided for one test over three seeds. After that test was
inconclusive for both candidates, a second look with three more seeds was
added for \acrshort{UMT}-L, the candidate with the larger gain, with its own rule
written before its runs at a stricter level (\cref{sec:statistics}). Before
those runs, it was also noted that the lower bound of the first look for
\acrshort{UMT}-L, $-0.0001$, lies within the Monte Carlo error of the bootstrap,
and the second look was therefore fixed at 100,000 resamples with
\acrshort{UMT}-L as its only candidate. The outcome of the first look as
computed was kept, and both looks are reported (\cref{sec:extensions}).

\paragraph{The decoder state.} Before its rule was written, the buffer was run
at test time on two holdout videos to check the code; those scores were not
used. The analysis of the remaining gap without the first window of each video
(\cref{sec:extensions}) was added after the scores by position had been seen,
and it is descriptive. That the primary comparison, with separately trained
models, could not have found equivalence at the spread between seeds was
noticed only after the result; the rule was applied unchanged.

\paragraph{Extensions that were not completed.} The predictability gate and
the sound-source maps stopped at the first stage of their rules, an outcome
that the rules provided for. The speaker maps were dropped after a feasibility
check that their rule did not foresee; the check used no gaze data, and no
model was trained with the maps, so no result of theirs was computed.

\paragraph{Uses of the validation set.} The DHF1K validation set is the test
set of this thesis. \Cref{tab:val-ledger} lists every use of it, including
the two decisions about the test-set submission that it informed.

\section{Commands for the Main Results}

The main DHF1K model with the \acrshort{ViT}-g encoder (seed 0) is trained, its blur is
tuned on the holdout, and it is scored once on the validation set with:

\begin{small}
\begin{verbatim}
python scripts/train.py --config configs/main_b1.yaml --set \
    run_dir=runs/vitg_s0 train.seed=0 train.epochs=3 \
    train.warmup_epochs=0.25 model.backbone=vjepa2_g "data.size=[256,256]"
python scripts/tune_sigma.py --run runs/vitg_s0 \
    --out runs/vitg/sigma_single.json
python scripts/evaluate.py --run runs/vitg_s0 --split val --metrics all \
    --sigma 8 --out runs/vitg/val_single_s0.json
\end{verbatim}
\end{small}

Passing \verb|--run| three times (seeds 0, 1 and 2) to the last two commands
evaluates the three-seed ensemble. The paired bootstrap against another
evaluation file is computed with \verb|python scripts/compare.py A.json B.json|.

The causal variant (\cref{sec:causal}) is trained and scored on every tenth
frame with:

\begin{small}
\begin{verbatim}
python scripts/train.py --config configs/main_b1.yaml --set \
    run_dir=runs/causal_g_s0 train.seed=0 train.epochs=3 \
    train.warmup_epochs=0.25 model.backbone=vjepa2_g \
    "data.size=[256,256]" data.target=last eval.frame_step=10
python scripts/evaluate.py --run runs/causal_g_s0 --split val \
    --metrics all --sigma 8 --frame-step 10 \
    --out runs/review/causal_g_val_s0.json
\end{verbatim}
\end{small}

\noindent \verb|scripts/evaluate.py| and \verb|scripts/tune_sigma.py| run a model
causally whenever its configuration names this target. The \acrshort{ViT}-L variant is
trained with the same command, in which \verb|vjepa2_l| replaces \verb|vjepa2_g|.
For the comparison of \cref{sec:main-results}, SalFoM's saved maps
are blurred by passing \verb|--sigma| together with \verb|--maps|.

The analyses of the second follow-up (the learning curve of
\cref{sec:learning-curve}, the prediction-error maps and their incremental
information in \cref{sec:g2}, and the measured times of \cref{sec:cost}) are
run with:

\begin{small}
\begin{verbatim}
bash scripts/queue_review_d.sh      # six training runs and their analysis
python scripts/g2_maps.py --out runs/g2/maps
python scripts/g2_incremental.py --out runs/review/g2_incremental.json
python -I scripts/time_inference.py --out runs/review/timing_run3.json
\end{verbatim}
\end{small}

\noindent The last command needs SalFoM's released code and weights, which are
not part of the repository of the thesis.

The extensions of \cref{sec:extensions} are run, in the order of the research
log, with the queues below; each writes the per-video holdout scores that the
tests read. The second-encoder tests, the decision on the sound-source maps and
the tests of the decoder state then follow from:

\begin{small}
\begin{verbatim}
bash scripts/queue_14a.sh && bash scripts/queue_14a2.sh   # seeds 0-2
bash scripts/queue_14d.sh                                 # seeds 3-5
bash scripts/queue_14e.sh                                 # gate
python -I scripts/ssl_maps.py --ckpt <EZ-VSL checkpoint> --rate 44100
bash scripts/queue_14f.sh                                 # sound maps
bash scripts/queue_16p1.sh                                # decoder state
python scripts/ablation_seeds.py --ref vitg_s0,vitg_s1,vitg_s2 \
    --variant dinov2_l=xp_dinov2_l_s0,xp_dinov2_l_s1,xp_dinov2_l_s2 \
    --variant umt_l=xp_umt_l_s0,xp_umt_l_s1,xp_umt_l_s2 \
    --out runs/review/expert_seeds_14a.json
R=vitg_s0,vitg_s1,vitg_s2,vitg_s3,vitg_s4,vitg_s5
U=xp_umt_l_s0,xp_umt_l_s1,xp_umt_l_s2
U=$U,xp_umt_l_s3,xp_umt_l_s4,xp_umt_l_s5
python scripts/ablation_seeds.py --ref $R --variant umt_l=$U \
    --level 97.5 --n 100000 --seed 0 --out runs/review/expert_seeds_14d.json
python scripts/summarize_ssl.py --stage 1
python scripts/inv_summary.py --part p1
python -I scripts/p1_first_window.py .
\end{verbatim}
\end{small}

\noindent The EZ-VSL checkpoint is the one its authors released for
VGG-Sound; it is read with tensor-only loading and is not part of the
repository.

\section{Forward Pass of the Decoder}

The listing below is the forward method of the decoder in
\texttt{surge/model.py}, which implements \cref{eq:fpn,eq:decoder,eq:output}
(the softmax over pixels is applied in the loss and at inference). The three
branches used only by the second encoder, the audio fusion and the modulation
of \cref{eq:modulation} are abbreviated; none of them is part of the final
model.

\begin{footnotesize}
\begin{verbatim}
def forward(self, taps, surprise=None, T_out=None, audio=None,
            expert_taps=None):
    p4, p8, p16a, p16b = self.fpn(taps)            # SimpleFPN-3D
    if self.expert_dim:                             # second encoder only
        ...
    if self.audio:                                  # audio experiment only
        ...
    if self.use_scm:                                # gate and sound maps only
        ...
    d16 = self.r16(self.fuse16(torch.cat([p16a, p16b], 1)))
    d8 = self.r8(p8 + F.interpolate(d16, size=p8.shape[2:], mode="nearest"))
    d4 = self.r4(p4 + F.interpolate(d8, size=p4.shape[2:], mode="nearest"))
    out = {"logits": self.head(d4)[:, 0]}           # B, T', 64, 64
    if self.aux and self.training:
        out["aux16"] = self.aux16(d16)[:, 0]
        out["aux8"] = self.aux8(d8)[:, 0]
    for k in list(out):
        x = out[k]
        if k != "logits":                           # aux maps to 64 x 64
            x = F.interpolate(x[:, None], size=out["logits"].shape[1:],
                              mode="trilinear", align_corners=False)[:, 0]
        if T_out is not None and x.shape[1] != T_out:   # tubelets -> frames
            x = F.interpolate(x[:, None], size=(T_out, *x.shape[2:]),
                              mode="trilinear", align_corners=False)[:, 0]
        if self.prior is not None:                  # + log G_theta
            x = x + self.prior(x.shape[-2], x.shape[-1], x.device)
        out[k] = x
    return out
\end{verbatim}
\end{footnotesize}
